% !TEX root = task-2-10-1.tex
% Editable task bank. Append new BankTask blocks at the end of the document.
\documentclass[12pt,a4paper]{article}
\usepackage{physics-textbook}
\usepackage{graphicx,needspace}
\usetikzlibrary{decorations.pathmorphing}
\newif\ifShowSolutions
\ifdefined\TutorEdition
  \ShowSolutionstrue
\else
  \ShowSolutionsfalse % Change to \ShowSolutionstrue to display solutions.
\fi
\newcounter{banktask}
\newsavebox{\BankTaskBox}
\newcommand{\BankTaskContent}[3]{%
  \noindent\textbf{\thebanktask.}\quad #2\par
  \ifShowSolutions
    \Needspace{50mm}% Keep the solution heading with its opening text or diagram.
    #3\par
    \ifdefined\TutorEdition
      \ifcsname TutorSource#1\endcsname
        \nopagebreak\smallskip{\small\csname TutorSource#1\endcsname\par}%
      \fi
    \fi
  \fi
}
\newcommand{\BankTask}[3]{%
  \refstepcounter{banktask}\par\medskip
  \sbox{\BankTaskBox}{\begin{minipage}{\linewidth}%
    \BankTaskContent{#1}{#2}{#3}\end{minipage}}%
  \ifdim\dimexpr\ht\BankTaskBox+\dp\BankTaskBox\relax>\dimexpr\textheight-20mm\relax
    \Needspace{35mm}\BankTaskContent{#1}{#2}{#3}%
  \else
    \Needspace{\dimexpr\ht\BankTaskBox+\dp\BankTaskBox\relax}%
    \noindent\usebox{\BankTaskBox}\par
  \fi
  \medskip
}
\begin{document}
\BookSubsection{2.10.1}{Плавление и кристаллизация}


% Source 70D177: https://ege.fipi.ru/bank/questions.php?search=1&proj=BA1F39653304A5B041B656915DC36B38&qid=70D177
\BankTask{70D177}{
В калориметр со льдом, имеющим температуру \(0\,^{\circ}\mathrm{C}\), добавили воду при комнатной температуре. Как изменятся в результате установления теплового равновесия удельная теплоёмкость воды и масса льда в калориметре?

Для каждой величины определите соответствующий характер изменения:
\begin{enumerate}
\item увеличится
\item уменьшится
\item не изменится
\end{enumerate}

\TaskChangeTable{Удельная теплоёмкость воды}{Масса льда}
}{
\smallskip\noindent\textbf{Решение от GPT}\par
Удельная теплоёмкость воды в принятой модели постоянна: \(c=4200\,\text{Дж}/(\text{кг}\cdot\text{К})\), поэтому не изменится (3). Вода при охлаждении отдаёт теплоту, которая расходуется на плавление льда: \(Q=\lambda\Delta m\). Масса льда уменьшится (2).

\(\boxed{32}\).
}

\end{document}
